\documentclass[11pt,a4paper]{article}
\usepackage[T1]{fontenc}
\usepackage{lmodern}
\usepackage[margin=26mm]{geometry}
\usepackage{amsmath,amssymb,amsthm,booktabs,microtype}
\usepackage[colorlinks=true,linkcolor=blue,urlcolor=blue,citecolor=blue]{hyperref}
\hypersetup{pdftitle={A certified median threshold for six-variable random 2-SAT},pdfauthor={Nefedov Sergey},pdfsubject={P2588: exact coefficient evaluation and independent verification}}
\usepackage{enumitem}
\setlist{nosep}
\newtheorem{theorem}{Theorem}
\newtheorem{lemma}[theorem]{Lemma}
\newcommand{\coeff}[1]{\left[#1\right]}
\title{A certified median threshold for\newline six-variable random 2-SAT}
\author{Nefedov Sergey\\\small Moscow State University, Mathematics Department}
\date{20 October 2026\\\small Version 3: reproducibility repository}
\begin{document}
\maketitle
\begin{abstract}
We evaluate the number of satisfiable thirteen-clause sets on six labeled Boolean
variables in the distinct-clause model. The exact count is
$2\,381\,437\,338\,960$, giving probability
$9\,922\,655\,579/21\,528\,597\,615<1/2$. Together with the known twelve-clause
count and monotonicity, this proves that thirteen is the first strict half-probability
crossing. The evaluation uses the existing enumeration theorem of Dovgal,
de Panafieu and Ravelomanana and is independently reproduced by a direct
assignment-mask dynamic program. We supply integer-polynomial certificates,
source code and executed verification logs. The contribution is a certified finite
evaluation of existing theory; no new general enumeration theorem is claimed.
The exact certificate arithmetic, the finite clause model and probability
monotonicity have also been checked in Lean~4. The identification of the
certificate coefficients with the satisfiable-formula counts is not yet
formalized, so the Lean verification is partial.
\end{abstract}

\section{Model and result}
Let $\mathcal C_n$ contain the four signed disjunctions on every unordered pair of
distinct variables from $x_0,\ldots,x_{n-1}$. Thus
$|\mathcal C_n|=4\binom n2=2n(n-1)$. A formula is an unordered subset of
$\mathcal C_n$. Variables are labeled; clauses cannot repeat. Write
\[
a_{n,m}=\#\{F\subseteq\mathcal C_n: |F|=m,\ F\text{ satisfiable}\},\qquad
P_{n,m}=\frac{a_{n,m}}{\binom{2n(n-1)}m}.
\]
The finite request P2588 of TheoremDB asks for the first $m$ with $P_{6,m}<1/2$.

\begin{theorem}
In this model the first such value is $m=13$. More precisely,
\[
a_{6,12}=805\,717\,285\,720,\qquad
a_{6,13}=2\,381\,437\,338\,960,
\]
and
\[
P_{6,13}=\frac{9\,922\,655\,579}{21\,528\,597\,615}
=0.46090580336205517\ldots.
\]
\end{theorem}

The twelve-clause value is in the published table of \cite{DPR}. We reproduced it
before evaluating the next coefficient. The exact integer comparisons are
\begin{center}
\begin{tabular}{r r r}
\toprule
$m$ & $\binom{60}m$ & $2a_{6,m}-\binom{60}m$\\
\midrule
12 & $1\,399\,358\,844\,975$ & $+212\,075\,726\,465$\\
13 & $5\,166\,863\,427\,600$ & $-403\,988\,749\,680$\\
\bottomrule
\end{tabular}
\end{center}

\begin{lemma}
For a fixed finite clause universe of size $C$, satisfiability probabilities of
uniform clause sets satisfy $P_{m+1}\le P_m$ for $0\le m<C$.
\end{lemma}
\begin{proof}
Choose a uniform $(m+1)$-set and delete a uniformly chosen clause. A fixed $m$-set
has $C-m$ extensions, and its probability under deletion is
\[
\frac{C-m}{(m+1)\binom C{m+1}}=\frac1{\binom C m}.
\]
The resulting $m$-set is therefore uniform. Deletion preserves satisfiability,
so the satisfiability indicator cannot decrease under this coupling. Taking
expectations proves the claim.
\end{proof}
Consequently the two integer comparisons prove the claimed first crossing,
once the coefficient evaluations below are certified.

\section{Exact coefficient evaluation}
Set $q=1+w$ and use exponential generating functions in $z$:
\[
G(z,w)=\sum_{r\ge0}q^{\binom r2}\frac{z^r}{r!},\qquad
H=G^{-1},\qquad E=G\odot_z H,\qquad A=\sqrt E,
\]
where $A(0,w)=1$ and $\odot_z$ is the exponential Hadamard product. Theorem~4.7
of \cite{DPR} gives, writing $A_r=r!\coeff{z^r}A$,
\begin{equation}\label{eq:count}
a_{n,m}=\coeff{w^m}\,2^n\sum_{k=0}^n\binom nk
q^{2(n-1)k-\binom k2}A_{n-k}(w).
\end{equation}
This is the source enumeration theorem, not a new formula of this note.
The normalization distinguishes labeled variables, unordered clause sets and
exponential coefficients explicitly.

For $H_r=r!\coeff{z^r}H$ and $E_r=q^{\binom r2}H_r$, formal inversion and
square-root extraction give
\begin{align}
H_0&=1,& H_r&=-\sum_{j=1}^r\binom rj q^{\binom j2}H_{r-j},\label{eq:inv}\\
A_0&=1,&2A_r&=E_r-\sum_{j=1}^{r-1}\binom rj A_jA_{r-j}.\label{eq:sqrt}
\end{align}
The primary implementation evaluates these identities as full polynomials in
$w$ over $\mathbb Q$, with exact rational arithmetic. Coefficients above $z^n$
cannot affect the requested result: each recurrence for degree $r$ uses only
degrees at most $r$. All $w$-degrees were retained, so no truncation in the clause
variable is involved. Full rows were computed for $n\le7$.

\subsection{An integer certificate}
A separately written evaluator uses ordinary powers of $q$ and the scaled
coefficients $R_r=2^r A_r$. Then
\begin{equation}\label{eq:scaled}
R_0=1,\qquad
2R_r=2^r E_r-\sum_{j=1}^{r-1}\binom rj R_jR_{r-j}.
\end{equation}
The final polynomial and its coefficient transform are
\begin{align}
S_n(q)&=\sum_{k=0}^n2^k\binom nk
q^{2(n-1)k-\binom k2}R_{n-k}(q),\\
a_{n,m}&=\sum_{j\ge m}\coeff{q^j}S_n(q)\binom jm.\label{eq:transform}
\end{align}
The stored certificate contains the integer coefficient arrays $H_r,E_r,R_r,S_r$.
The stand-alone checker verifies $GH=1$, the Hadamard identity and
$R(z)^2=E(2z)$, then checks \eqref{eq:count}--\eqref{eq:transform} and all saved
coefficient rows. It imports neither evaluator and needs no symbolic algebra
package. Uniqueness of the formal inverse and the square root with constant
coefficient~1 justifies the finite identity certificate.

\section{Independent direct enumeration}
The second counting method uses no generating functions. Encode the $64$ truth
assignments by bits of a 64-bit unsigned integer. For a clause $c$, let $M(c)$
have a bit precisely at assignments satisfying $c$. Each clause mask has $48$
bits. The solution set of a formula is exactly $\bigcap_{c\in F}M(c)$.

After the first $i$ clauses, let $D_i(M,k)$ count $k$-element subsets with exact
solution mask $M$. Initially $D_0(\mathrm{ALL},0)=1$. Every clause produces the
two transitions
\[
(M,k)\longrightarrow(M,k),\qquad
(M,k)\longrightarrow(M\mathbin{\&}M(c),k+1).
\]
The two alternatives partition all subsets by absence or presence of the new
clause. Induction therefore proves that the DP counts each labeled clause set
exactly once. Equal masks are merged by addition of multiplicities, not by
identification of formulas. An empty mask can be discarded permanently;
degrees above~13 can also be discarded.

\paragraph{In-place evaluation.}
To reduce memory, take a snapshot of the existing state nodes before each clause
and process them in increasing mask population count. A changed mask is a proper
subset and has strictly smaller population count. Its node is therefore already
processed or is new and absent from the snapshot. Hence no earlier transition
can alter the source polynomial before it is read. If the mask is fixed, update
degrees in decreasing order, multiplying its polynomial by $1+w$. This implements
the same recurrence as the simple two-buffer reference program.

For $k\le13$, every count is bounded by $\binom{60}{13}<2^{64}$, so the counting
arithmetic cannot overflow. The state checksums intentionally use arithmetic
modulo $2^{64}$ and play no role in the counts.

\paragraph{Executed output.}
All $60$ steps completed. There were $24\,444\,343$ nonempty final masks; the last
step was recorded at $117.657$ seconds. All fourteen counts, $0\le m\le13$,
agreed with the primary calculation. The final two entries were exactly the
values in Theorem~1. The run logs retain a deterministic order-independent
64-bit checksum for every state table. SHA-256 digests cover distributed source,
certificates, logs and outputs. The large state tables are regenerated on demand
and are not distributed; the 64-bit checksums are not cryptographic certificates.

\subsection{Additional executed checks}
\begin{itemize}
\item Every one of the $2^{24}=16\,777\,216$ clause sets at $n=4$ was enumerated
separately. The complete row agrees in all $25$ degrees. A simple two-buffer
Python DP also agrees for every degree at $n\le4$.
\item For all $4096$ formulas at $n=3$, truth-table satisfiability agrees with
the implication-graph SCC criterion. A further fixed-seed test checks $1220$
formulas at $n=6$. These tests are not probability estimates.
\item Both unmodified coefficient functions from the pinned upstream code were
executed at $(6,12)$ and $(6,13)$ and returned the same counts.
\item The published totals for $n\le7$ were reproduced. In particular,
\[\sum_m a_{6,m}=2\,034\,602\,766\,692\,687.\]
Its complement is
$1\,150\,886\,901\,840\,154\,289$, and the two sum to $2^{60}$.
\end{itemize}

\section{Reproducibility and scope}
The intended repository URL for these materials is
\begin{center}
\url{https://github.com/Sergpreneur/P2588-2SAT-Threshold}.
\end{center}
Public upload is pending for this version; the accompanying archive contains
the complete materials prepared for that repository.
The package includes both polynomial evaluators, the independent C++ mask DP,
the exhaustive checker, full data, certificates, logs, version information and
the Russian report. From the package root:
\begin{verbatim}
python3 verify/verify_count.py
python3 -m unittest discover -s tests -v
python3 verify/verify_count.py --rerun-dp
\end{verbatim}
The first command verifies the compact certificate and saved outputs; the last
also recompiles and reexecutes the independent large DP. The recorded
environment used Python~3.12.14 and g++~13.3.0 with \texttt{-O3 -std=c++17}.
SymPy~1.14.0 was used only to run upstream code. The main programs and the compact
checker use the Python standard library. The C++ DP requires several GiB of
memory. These computational checks are complemented by the partial Lean
verification described below.

\section{Lean 4 verification and its scope}\label{sec:lean}
The accompanying formal project was compiled with Lean~4.19.0 and mathlib
v4.19.0. The pinned mathlib revision is
\begin{center}\texttt{c44e0c8ee63ca166450922a373c7409c5d26b00b}.\end{center}
The standard
\texttt{lake build} completed successfully; the build and axiom reports are
included in the \texttt{lean/} directory, together with the formal sources.

The arithmetic modules check the finite integer-polynomial identities for
$H_r,E_r,R_r$ through degree six, the assembly of $S_6$, the binomial transforms
at twelve and thirteen clauses, both binomial denominators, the exact signed
gaps, and the reduced probability. They also check every strict above-half
comparison for the coefficients at $0\le m<13$. These modules depend only on
Lean's standard library and use ordinary \texttt{decide}, whose proof terms
are checked by the kernel. The reported arithmetic theorems have no axiom
dependencies; neither \texttt{native\_decide} nor an external computational
oracle is used.

The separate model module defines a clause by a variable pair $i<j$ and two
Boolean signs, represents formulas as finite sets, and defines satisfiability
by an existential truth assignment. Lean proves that the clause universe has
cardinality sixty and that there are $\binom{60}{m}$ formulas of size $m$.
It also proves deletion preserves satisfiability and that the corresponding
probability is nonincreasing. This last proof uses the local
Lubell--Yamamoto--Meshalkin inequality from mathlib. Its axiom dependencies are
the standard \texttt{propext}, \texttt{Classical.choice} and \texttt{Quot.sound}.
There are no admitted proof terms or new enumeration axioms.

The remaining formal gap is explicit. Writing \texttt{satCount} for the count
of satisfiable formulas in the Lean model, the final theorem
\texttt{threshold13\_of\_enumeration\_bridge} assumes
\[
\begin{split}
\texttt{satCount}(12)&=\coeff{w^{12}}S_6(1+w),\\
\texttt{satCount}(13)&=\coeff{w^{13}}S_6(1+w).
\end{split}
\]
Under these hypotheses it proves that thirteen is the first strict crossing.
The hypotheses are not proved inside Lean: the specialization of the source
enumeration theorem, or a verified counting procedure connecting the coefficient
certificate to this finite model, remains to be formalized. Thus the numerical
and monotonicity parts are kernel-checked, while the complete counting argument
and the unconditional threshold theorem are not yet fully formalized in Lean.
The mathematical result rests on \cite{DPR} and the independent executed
enumeration described above.

\paragraph{Literature and priority.}
The general counting problem is already solved by \cite{DPR}. The latest
displayed arXiv version was v2, and the inspected GitLab master head remained
commit \texttt{346079fa8606ebd48f4557bcef7bc1f73a2aed4d}. Its notebook displays
the row only through $m=12$, but also computes totals by summing all degrees.
Thus it already evaluated $a_{6,13}$ internally. A bounded search found no separate
printed value for this coefficient, which does not establish priority. This note
provides an explicit certified evaluation and closes the finite condition in
P2588; it does not claim a new general theorem or a new asymptotic 2-SAT threshold.
The primary journal article is article~7, correcting article~8 in the supplied
problem metadata.

\paragraph{AI assistance.}
ChatGPT (OpenAI) was used to derive and implement the exact recurrences, design the
independent verification, execute computations, audit sources and draft this
note. The mathematical argument and computational claims should be assessed
from the supplied code and certificates. The exact scope of the Lean checks
is stated in Section~\ref{sec:lean}. This version records no journal submission.

\begin{thebibliography}{9}
\bibitem{DPR} S. Dovgal, E. de Panafieu and V. Ravelomanana.
\emph{Exact enumeration of satisfiable 2-SAT formulae}.
Combinatorial Theory \textbf{3}(2) (2023), article~7.
\href{https://doi.org/10.5070/C63261985}{doi:10.5070/C63261985}.
\href{https://arxiv.org/abs/2108.08067v2}{arXiv:2108.08067v2}.
\bibitem{code} S. Dovgal and E. de Panafieu. Companion repository,
\emph{Enumeration-2sat-aux}, commit
\texttt{346079fa8606ebd48f4557bcef7bc1f73a2aed4d}.
\url{https://gitlab.com/sergey-dovgal/enumeration-2sat-aux}.
\bibitem{problem} TheoremDB contributors. P2588,
\emph{Median satisfiability threshold for a six-variable clause set}.
\url{https://theoremdb.org/statements/P2588/}. Accessed 20 October 2026.
\end{thebibliography}
\end{document}
